\documentclass[11pt,a4paper]{ctexart}
\usepackage[margin=2.4cm]{geometry}
\usepackage{amsmath,amssymb,booktabs,hyperref,xcolor}
\hypersetup{colorlinks=true,linkcolor=blue,urlcolor=blue}

\title{纽约出租车区域推荐：有限时域规划与可复现实证}
\author{Zefan Cai}
\date{2026 年 7 月}

\begin{document}
\maketitle

\begin{abstract}
本文研究给定司机当前位置和时刻时的 Top-3 出租车候客区域推荐。系统使用纽约市 TLC 2023 年 1 月黄出租车订单，构建按星期、半小时和区域聚合的需求与车费统计、有向 OD 行程时间图、乘客下车区域转移模型，并比较热门区域、单步效用和有限时域规划。本文只报告由机器可读实验产物复现的结果，并明确区分参考效用排序指标、固定单司机模拟器收益和真实市场因果收益。公开静态诊断中，两步策略的 NDCG@3 为 0.9565、Hit@3 为 0.9714；100 个配对随机种子的七日模拟中，两步策略平均日基础车费为 570.61 美元，相对单步策略增加 21.84 美元，95\% bootstrap 区间为 [5.00,39.53]。由于模拟器不包含司机竞争、需求耗尽、拥堵和供需反馈，结果不能解释为生产部署收益。
\end{abstract}

\section{数据与时间边界}

训练期为 2023-01-01 至 2023-01-24，公开验证期为 2023-01-25 至 2023-01-31。训练订单最晚结束于 1 月 24 日 23:59:56，验证订单最早开始于 1 月 25 日 00:00:00。审计未发现跨集合精确重复订单，区域时段统计可以从清洗训练集逐项重建。

清洗规则包括日期边界、必要字段、区域编号、车费、时长、距离、速度和重复订单过滤。当前代码可直接从官方月度 parquet 自动生成按时间切分的未清洗训练与验证输入。

\section{方法}

设 $D(s,z)$ 为训练期状态 $s$ 和区域 $z$ 的订单数，$\bar f(s,z)$ 为平均基础车费。接单概率是场景假设
\[
p(s,z)=\frac{D(s,z)}{D(s,z)+240}.
\]
一步终端价值为
\[
V_1(s,z)=p(s,z)\bar f(s,z).
\]
对于 $h>1$，有限时域续接价值为
\[
V_h(s,z)=p(s,z)\left[\bar f(s,z)+\gamma\sum_{z'}P(z'\mid z)V_{h-1}(s+1+\tau_z,z')\right]
+(1-p(s,z))\gamma V_{h-1}(s+1,z).
\]
初始空驶动作再按离散移动时段归一化。未来续接策略固定为在到达区域等待，因此该方法属于截断 lookahead，而不是完整的两步 Bellman 最优规划。

仓库另提供修正后的模型式 MDP：动作包含空驶，成功后按 OD 分布到达下车区，失败后留在候选区，两种结果均推进时间，并使用同步 Bellman backup。

\section{可复现实验}

\subsection{静态参考目标诊断}

\begin{table}[h]
\centering
\begin{tabular}{lrrr}
\toprule
策略 & NDCG@3 & Hit@3 & 参考效用@1 \\
\midrule
热门区域 & 0.7846 & 0.5842 & 19.4299 \\
单步效用 & 0.9024 & 0.8804 & 25.0589 \\
两步规划 & \textbf{0.9565} & \textbf{0.9714} & \textbf{27.5855} \\
\bottomrule
\end{tabular}
\end{table}

这些指标比较策略与公开两步参考效用的排序一致性，不是历史司机行为、真实最优动作或反事实收益。

\subsection{固定模拟器}

\begin{table}[h]
\centering
\begin{tabular}{lrrr}
\toprule
策略 & 平均日车费 & 标准差 & 平均接单数 \\
\midrule
热门区域 & 431.21 & 37.92 & 133.91 \\
单步效用 & 548.77 & 72.62 & 107.01 \\
两步规划 & \textbf{570.61} & 65.23 & 91.69 \\
\bottomrule
\end{tabular}
\end{table}

两步与单步使用相同随机种子配对比较：均值差为 21.84 美元/日，bootstrap 95\% 区间 [5.00,39.53]，配对 $t$ 检验 $p=0.0151$，Wilcoxon $p=0.000408$，Cohen's $d_z=0.247$。这是小效应，且只反映固定模拟器内的 Monte Carlo 波动。

\subsection{规划时域}

\begin{table}[h]
\centering
\begin{tabular}{lrr}
\toprule
时域 & NDCG@3 & 平均日车费 \\
\midrule
1 & \textbf{0.9582} & 569.78 \\
2 & 0.9565 & 570.61 \\
3 & 0.9549 & 573.47 \\
5 & 0.9525 & \textbf{575.97} \\
自适应 & 0.9559 & 573.31 \\
\bottomrule
\end{tabular}
\end{table}

参考 NDCG 与模拟车费并不单调一致，因此不能用高 NDCG 直接推断更高收入。

\section{鲁棒性与市场集中}

随机遮蔽 10\% 区域时段统计后，NDCG 从 0.9565 降至 0.9042，Hit@3 从 0.9714 降至 0.7872，说明缺失数据是主要不稳定来源。

两步策略的加权机场曝光为 70.33\%，其中 JFK 单区约 55.0\%；曝光 Gini 为 0.982，有效曝光区域数仅 5.51。单司机模拟器未建模机场排队和推荐饱和，因此这一集中度构成显著部署风险。

\section{反事实评估与强化学习边界}

TLC 订单不包含展示给司机的空驶推荐、行为策略概率或司机接受行为，因此无法从该表合法识别 IPS、SNIPS、Doubly Robust、CQL 或 BCQ 的真实推荐策略结果。仓库中的 Q-learning 是在历史统计构造的模拟器内进行在线探索训练，不属于 logged-data offline RL。

\section{模拟器限制}

公开 rollout 只模拟一个司机，历史需求单元不可耗尽，且不包含司机竞争、拥堵、供需反馈、机场队列和市场均衡。因而本文结果只支持同一模拟器内的策略比较，不支持真实部署增收结论。

\section{可复现性}

完整指标快照保存在 \texttt{outputs/reference\_metrics.json}，Markdown 报告由 \texttt{scripts/generate\_evaluation\_report.py} 自动生成。参数选择脚本真实运行每个配置，不再返回硬编码指标。完整研究有效性审计见 \texttt{outputs/research\_grade\_audit.md}。

\section{结论}

两步规划在给定参考目标和固定单司机模拟器中优于两个基础方法，但改进幅度、科学新颖性和外部有效性均应谨慎陈述。后续工作需要多月滚动测试、供给数据、多司机需求耗尽模拟器以及随机化线上日志。

\end{document}
